% =====================================================================
\chapter*{Instrucciones}
\addcontentsline{toc}{chapter}{Instrucciones}
\markboth{Instrucciones}{}
% =====================================================================

Este examen cierra la Fase~I del curso. Evalúa las herramientas matemáticas que se emplean en el resto del programa: la probabilidad del siguiente token, los vectores y las matrices, los gradientes, la estimación estadística y la medida de la información en bits.

\section*{Condiciones del examen}
\begin{center}
\begin{tabularx}{\textwidth}{@{}lX@{}}
\toprule
Duración & 2 horas. Reparto orientativo: 30 minutos para la parte~A y 90 minutos para la parte~B. \\
Material permitido & Calculadora científica y una hoja A4 de fórmulas escrita a mano por el estudiante. No se permiten apuntes, libros ni dispositivos con conexión. \\
Puntuación total & 100 puntos: parte~A, 30 puntos; parte~B, 70 puntos. \\
Aprobado & 70 puntos o más (70\,\%). \\
Módulos evaluados & M01 a M06 (véase la tabla siguiente). \\
\bottomrule
\end{tabularx}
\end{center}

\section*{Normas}
\begin{enumerate}[label=\arabic*.]
  \item \textbf{Parte A. Cuestiones} (30 puntos): 20 preguntas de opción múltiple, de 1,5 puntos cada una. Cada pregunta tiene una sola respuesta correcta. Las respuestas en blanco y las erróneas valen 0 puntos; las erróneas no restan.
  \item \textbf{Parte B. Problemas} (70 puntos): 6 problemas de desarrollo, uno por módulo. La puntuación de cada apartado figura entre corchetes.
  \item En la parte~B debe escribirse el razonamiento completo. Un resultado correcto sin justificación se puntúa con la mitad de los puntos del apartado.
  \item Los resultados se dan con 3 cifras significativas salvo que se indique otra cosa; en los pasos intermedios conviene conservar 4 o más. Se escribe $\ln$ para el logaritmo natural (resultados en nats) y $\log_2$ para el logaritmo en base~2 (resultados en bits).
  \item El \emph{proyecto de fase} (capítulo final del enunciado) se entrega aparte, en el plazo de una semana, y cuenta en la nota de la fase, no en la de este examen.
\end{enumerate}

\section*{Módulos evaluados}
\begin{center}
\begin{tabular}{@{}cllc@{}}
\toprule
Módulo & Título & Cuestiones & Problema \\
\midrule
M01 & Modelos de lenguaje probabilísticos & 1--4 & B.1 \\
M02 & Álgebra lineal I & 5--6 & B.2 \\
M03 & Álgebra lineal II & 7--11 & B.3 \\
M04 & Cálculo para el aprendizaje automático & 12--14 & B.4 \\
M05 & Probabilidad y estadística & 15--17 & B.5 \\
M06 & Teoría de la información & 18--20 & B.6 \\
\bottomrule
\end{tabular}
\end{center}

\section*{Procedencia de los problemas}
Los problemas B.3, B.4, B.5 y B.6 están adaptados, total o parcialmente, de exámenes y hojas de problemas de cursos universitarios (MIT 18.06, Stanford CS224N, MIT 18.05 y la asignatura \emph{Information Theory} de la Universidad de Cambridge); la fuente se indica al pie de cada enunciado y en la bibliografía. Los enunciados son de redacción propia y los datos se han cambiado para que puedan resolverse a mano. Las soluciones completas están en el \cref{ap:soluciones}, y la tabla de autoevaluación, en el \cref{ap:autoevaluacion}.
